\section{Types}
\label{sec:types}

The type structure underlying \algst comprises three parts: the kind
structure, the type language and protocol definitions. The type
language and protocol definitions are mutually recursive.  To
simplify the exposition, the formalized language is based on linear
System~F. While we formalize protocol
definitions, we do not formalize the standard treatment of (linear)
algebraic datatypes and pattern matching, although we use it in examples.

The type language encompasses types with different features
(as illustrated in \cref{sec:algebr-sess-types}). 
For a correct type abstraction, we need to endow the system with a classification 
of types through kinds.
%
% \paragraph{Kinds}
% \label{sec:kinds}
% \label{sec:what-your-kind}
%
\algst distinguishes three kinds.
\begin{description}
\item[$\kinds$] classifies session types, which in turn classify
  endpoints of communication channels.
% \item Kind $\kindm$ classifies types whose serialization is hard coded. These are
%   primarily primitive types like $\TInt$ and $\TBool$. We also
%   include session types.
\item[$\kindt$] classifies all types that classify run-time
  values.
\item[$\kindp$] classifies protocol types. The kind $\kindp$ subsumes any other
  kind by an implicit lifting that associates a type with a protocol
  suitable to transmit it. A protocol type per se classifies pure
  behavior, but no run-time values.
\end{description}

Kinds are linearly ordered in a subkinding relation, which is the
reflexive transitive closure of the strict ordering $\kinds
< \kindt < \kindp$. Type formation includes a standard kind
subsumption rule.

% The subsumption $\kindt < \kindp$ allows us to write types like
% $\TIn{(\TFun[] TU)}S$ or $\TOut{(\TForall\alpha\kind T)}S$. Such
% protocols are currently empty because they cannot be implemented in
% \algst. Alternatively, we could define another kind, $\kindl$, say,
% such that $\kindl < \kindp$, $\kindl < \kindt$, and $\kindt \not<
% \kindp$, but we chose to keep kinding simple.

% \paragraph{Functional types, session types, and protocol types}
% \label{sec:session-types}

% TYPES

% \Cref{fig:grammar-of-types} contains the grammar of types.

Here is the grammar of types:
%
% \begin{figure}[tp]
\begin{align*}
  \typec S, \typec T, \typec U \grmeq&
  \TUnit \grmor \TFun[] T U \grmor \TPair{T}{U} \grmor
  \TForall \alpha \kind T \grmor \typec{\alpha} \grmor
  && \text{functional types}
  \\
  & \TIn TS \grmor \TOut TS \grmor \TEndW \grmor \TEndT \grmor \TDual S \grmor
  && \text{session types}
  % \\
  % & \typec{\alpha} \grmor
  % & \text{functional and session types}
  \\
  & \TPApp \proto {\overline{T}} \grmor \pneg T 
  && \text{protocol types}
\end{align*}
%   \caption{Grammar of types}
%   \label{fig:grammar-of-types}
% \end{figure}
%
We mostly use the metavariable $\typec S$ to range over session
types (kind $\kinds$). Metavariables $\typec T, \typec U$ range over functional types
(kind $\kindt$) and protocol types (kind $\kindp$).
Functional types comprise the usual types of linear System F: unit,
function, product, universal quantification, and type variables.
% The handling of the unit type serves as blueprint for other primitive types.
The session type operators are receiving, sending,
passive and active termination, and the dual operator. The dual
operator swaps the direction of all communications in the spine of a
session type. Thus, a session type is a
sequence of communications either terminated by $\TEndW$ or by
$\TEndT$, or extensible with a type variable $\typec\alpha$ of kind $\kinds$.


The protocol type
$\TPApp \proto {\overline{T}}$ refers to the protocol declared as
$\proto$ with parameters $\typec{\overline{T}}$. The
operator ``\typec{$-$}'' reverses the direction of communication of its
protocol parameter. While duality transforms session types
outside-in, the reverse operator transforms protocols inside-out.

%
% Unlike in previous work on context-free session types
% \cite{DBLP:conf/icfp/ThiemannV16,DBLP:journals/toplas/Padovani19},
% there is no sequencing operator $\TSeq\_\_$, but sending and receiving
% types have their traditional prefix form,

A kind context $\Delta$ maps type variables to kinds and protocol
names to first-order arrow kinds that determine the
number of parameters of the  protocol.
Formally, we should write $\kindas\varepsilon\kindp$ for the kind of a
protocol constructor without parameters, but we often abbreviate it to $\kindp$.
% In informal settings, we just write $\kind$.
\begin{align*}
  \Delta & \grmeq \Empty
           \grmor \Delta, \tbind\alpha\kind
           \grmor \Delta, \tbind\proto{\kinda\kindp\kindp}
  % \\
  % \Pi & ::= M
  %       \grmor \protocolsfull; \Pi
\end{align*}

% TYPE FORMATION

\begin{figure}[t!]
  \declrel{Type formation (synthesis, check-against)}{$\typeSynth T
    \kind$,\quad $\typeAgainst T \kind$}
  \begin{mathpar}
    % FUNCTIONAL TYPES
    \infrule{\rulenametypeUnit}{}{\typeSynth\TUnit\kindt}

    \infrule{\rulenametypeArrow}
    {\typeAgainst T\kindt \\ \typeAgainst U\kindt}
    {\typeSynth{\TFun[] TU}\kindt}
    
    \infrule{\rulenametypePair}
    {\typeAgainst T\kindt \\ \typeAgainst U\kindt}
    {\typeSynth{\TPair TU}\kindt}

    \infrule{\rulenametypeVar}{\tbind\alpha\kind \in \Delta}{\typeSynth\alpha\kind}
    
    \infrule{{\rulenametypePoly}}
    {\typeAgainst[\Delta,\tbind\alpha\kind] T {\kindc{\kindt}}}
    {\typeSynth {\TForall \alpha\kind T}{\kindt}}

    % SESSION TYPES
    \infrule{\rulenametypeDotIn}
    {\typeAgainst T \kindp \\ \typeAgainst S \kinds}
    {\typeSynth{?T.S} \kinds}

    \infrule{\rulenametypeDotOut}
    {\typeAgainst T \kindp \\ \typeAgainst S \kinds}
    {\typeSynth{!T.S} \kinds}

    \infrule{\rulenametypeEnd}{}{\typeSynth\TEnd\kinds}

    \infrule{\rulenametypeDualS}
    {\typeAgainst S \kinds}
    {\typeSynth{\TDual S} \kinds}
    
    % GUARDED
    \infrule{\rulenametypeProtocol}
    { \protocolsfull \\
      \varnothing\ne J\subseteq I \\
      \typeAgainst{\overline T}{\overline{\kindp}}}
    {\typeSynth{\TPAppSet \proto J {\overline{T}}}\kindp}

    \infrule{{\rulenametypeMsgNeg}}
    {\typeAgainst{T}{\kindp}}
    {\typeSynth{{-T}}\kindp}

    \infrule{\rulenametypeSub}
    {\typeSynth T\kind
      \\ \isSubKind \kind \kind'}
    {\typeAgainst T\kind'}
  \end{mathpar}
  % \declrel{Type formation (check-against)}{$\typeAgainst T \kind$}
  % \begin{mathpar}
  %   \infrule{\rulenametypeSub}
  %   {\typeSynth T\kind
  %     \\ \isSubKind \kind \kind'}
  %   {\typeAgainst T\kind'}
  % \end{mathpar}
  \caption{Algorithmic type formation rules}
  \label{fig:type-formation}
\end{figure}

% \infrule{\rulenametypeDot}
% {\typeAgainst G \kindg \\ \typeAgainst S \kinds}
% {\typeSynth{G.S} \kinds}
% 
% \infrule{\rulenametypeUp}{\typeAgainst T \kindt}{
% \typeSynth{\TUp T}\kindp}
% 
% \infrule{{\rulenametypeMsgPos}}
% {\typeAgainst{T}{\kindp}
% }
%   {\typeSynth{{+T}}\kindp}
%   
%   \infrule{\rulenametypeInt}{}{\typeSynth\TInt\kindt}
%   
%   \infrule{\rulenametypeBool}{}{\typeSynth\TBool\kindt}
%   

%%%   Local Variables:
%%%   mode: latex
%%%   TeX-master: "main"
%%%   End:

%\begin{figure}[t!]
	\declrel{Type formation}{$\hasType T\kind$}
\begin{mathpar}
  \infrule{\rulenametypeUnit}{}{\hasType\TUnit\kindt}
  
  \infrule{{\rulenametypeArrow}}{
    \hasType T \kindt
    \\ \hasType U \kindt
  }{
    \hasType{\TFun[] T U}{\kindt}}

  % \infrule{\rulenametypeInt}{}{\hasType\TInt\kindt}
  % \infrule{\rulenametypeBool}{}{\hasType\TBool\kindt}
  \infrule{\rulenametypePair}{\hasType T \kindt \\ \hasType U \kindt}{\hasType{\TPair TU}\kindt}

  \infrule{\rulenametypeVar}{\typec{\alpha}\colon\kind \in \Delta}{\hasType\alpha \kind}
  \\
  \infrule{{\rulenametypePoly}}{
    \hasType[\Delta,\typec{\alpha}\colon\kind]{\typec{T}}{\kindt}
  }{
    \hasType{\forall (\typec{\alpha}\colon\kind) \typec{T}}{\kindt}}
  
  \infrule{\rulenametypeEndW}{}{\hasType\TEndW{\kinds}}

  \infrule{\rulenametypeEndT}{}{\hasType\TEndT{\kinds}}
 
  \infrule{\rulenametypeMsg}{\hasType T \kindp \\
    \hasType S \kinds}{\hasType
    {! T.S}{\kinds} \\ \hasType {? T.S}{\kinds}}

  \infrule{\rulenametypeDualS}{\hasType S {\kinds}}{\hasType{\TDual S} {\kinds}}

  \infrule{\rulenametypeAssume}{
    \typec{\proto : \overline{\kindp} \to \kindp} \in \Delta \\
    \hasType{\overline T}{\overline {\kindp}}
  }{\hasType{\TPApp \proto {\overline{T}}}{\kindp}}
  
  \infrule{\rulenametypeMinus}{\hasType T \kindp} {\hasType {\pneg T} \kindp}
  
  \infrule{\rulenametypeSub}{\hasType{\typec T}\kind \\ \kind \le \kind'}{\hasType{\typec T}\kind'}
  % \\
  % \infrule{\rulenametypeUp}{\hasType T \kindt}{\hasType{\TUp T} \kindp}

  % \infrule{\rulenametypePlus}{\hasType T \kindp}{\hasType {\ppos T} \kindp}

  % \textcolor{red}{\infrule{\rulenametypeSemi}{
  %   \hasType T \kindp \\ \hasType U \kindp}{
  %   \hasType {T;U} \kindp
  % }}
  % \\
  % \\
  %   % \kind \in \{\kindt, \kindp\} \\
\end{mathpar}
  \caption{Type formation}
  \label{fig:type-formation}
\end{figure}

%%% Local Variables:
%%% mode: latex
%%% TeX-master: "main"
%%% End:


\Cref{fig:type-formation} contains the type formation rules. We
present them in algorithmic form from the beginning. The judgment
$\typeSynth T \kind$ assigns  to type $\typec T$ its minimal kind $\kind$ under kind
context $\Delta$.  The corresponding checking judgment $\typeAgainst T \kind$ takes
inputs $\typec T$ and $\kind$ and checks if the synthesized kind is a
subkind of the expected one.

The types for unit (\rulenametypeUnit),  functions
(\rulenametypeArrow), pairs (\rulenametypePair), and universal
polymorphism (\rulenametypePoly)
have kind $\kindt$ like their constituents. Variables have any kind
that has been assigned to them (\rulenametypeVar).
%
\algst's notion of session type distinguishes passive $\TEndW$ and active
$\TEndT$ termination of a session (\rulenametypeEndW, \rulenametypeEndT), and the
dual operator (\rulenametypeDual). Conventional session types require types for
the values exchanged in messages; here we ask for protocols instead
(\rulenametypeDotIn and \rulenametypeDotOut).
%
Rule {\rulenametypeProtocol} builds a protocol type from a protocol name
$\proto$. We assume that the kind context contains a binding for $\proto$ that
determines its arity. All parameters to a protocol must be protocols
themselves. The latter constitutes no restriction because any type can be
lifted to a protocol. % We omit empty parameter lists.
A protocol type can change direction using the reverse operator (rule {\rulenametypeMsgNeg}).

% \begin{lemma}[Subsumption] If $\typeAgainst T\kind$ and $\isSubKind \kind
%   \kind'$, then $\typeAgainst T{\kind'}$.
% \end{lemma}
% \begin{proof}
%   Immediate from \rulenametypeSub and transitivity of subkinding.
% \end{proof}

% \paragraph{Protocol definitions}
% \label{sec:protocol-definitions}
% \label{sec:decl-spec-prot}

Algebraic protocol types $\TPApp \proto {\overline{\alpha}}$ are defined analogously to algebraic
datatypes.
\begin{gather*}
  \protocolsfull
  % \qquad
  % \protocolsdata
\end{gather*}
This declaration defines the protocol type constructor $\proto$ that takes as
many parameters as indicated by $\typec{\overline\alpha}$. The protocol has choices
indexed by $i\in I$. Each choice is tagged by a \emph{selector tag} $\termc{\constr_i}$
that guards a sequence of subprotocols $\typec{\overline{T_i}}$, to be processed in that
order. We assume that selector tags are globally unique. Program may use any
full application $\TPApp \proto {\overline U}$ as a protocol type.
%
% For a data protocol, we can indicate our intent by using the $\datak$
% keyword instead of $\protocolk$. Moreover, all parameters must be
% proper types, not protocols.
% \begin{gather*}
%   \protocolsdata
% \end{gather*}
%
%
% We start with a declarative rule set to determine well-formedness of a
% protocol definition.
%
We check a single protocol declaration using the formation rules where $\Delta$ may
refer to previously defined protocols:\footnote{For mutually
  recursive protocols $\typec{\overline\proto}$, we add the
  assumptions for all $\typec{\overline\proto}$ when checking each individual $\typec{\proto_i}$.}
\begin{mathpar}
  \inferrule{
    \protocolsfull \\
    \typeAgainst[\Delta, \tbind\proto{\kinda\kindp\kindp},
    \tbind{\overline{\alpha}}{\overline{\kindp}}]
    {\overline{T_i}}{\kindp}\\
    (\forall i\in I)
  }{
    \typeSynth{\proto}{\kinda\kindp\kindp}
  }
  % \\
  % \inferrule{
  %   \protocolsdata \\
  %   \typeAgainst[\Delta, \tbind\data{\kinda\kindt\kindt},
  %   \tbind{\overline{\alpha}}{\overline{\kindt}}]
  %   {\overline{T_i}}{\kindt}\\
  %   (\forall i\in I)
  % }{
  %   \typeSynth{\data}{\kinda\kindt\kindt}
  % }
\end{mathpar}

% \paragraph{Type conversion}
% \label{sec:type-conversion}

% TYPE CONVERSION

\begin{figure}[t!]
  \declrel{Type conversion}{$\isConv T T \kind$}
  % 
  \begin{mathpar}
    % EQUIVALENCE
    \infrule{\rulenametconvReflex}{\typeAgainst T \kind}{\isConv T T \kind}

    \infrule{\rulenametconvSymm}{\isConv {T_1} {T_2} \kind}{\isConv {T_2} {T_1} \kind}

    \infrule{\rulenametconvTrans}{\isConv {T_1} {T_2} \kind \\ \isConv {T_2} {T_3} \kind}{\isConv {T_1} {T_3} \kind}

    \infrule{\rulenametconvDualEndW}{}{\isConv{\TDual\TEndW}\TEndT\kinds}
    
    \infrule{\rulenametconvDualEndT}{}{\isConv{\TDual\TEndT}\TEndW\kinds}
    
    \infrule{\rulenametconvDualInM}{\typeAgainst T \kindp \\ \typeAgainst S {\kinds}}{
      \isConv{\TDual{(\TIn TS)}} {\TOut T{(\TDual S)}} \kinds}

    \infrule{\rulenametconvDualOutM}{\typeAgainst T \kindp \\ \typeAgainst S {\kinds}}{
      \isConv{\TDual{(\TOut TS)}} {\TIn T{(\TDual S)}} \kinds}

    \infrule{\rulenametconvDoubleS}{\typeAgainst S {\kinds}}{\isConv{\TDual{(\TDual S)}} S {\kinds}}

    \infrule{\rulenametconvSub}{
      \isConv {T_1}{T_2} {\kind} \\ \isSubKind\kind{\kind'}
    }{
      \isConv {T_1}{T_2} {\kind'}
    }
    
    \infrule{\rulenametconvNegIn}{\typeAgainst T \kindp \\ \typeAgainst S \kinds}{\isConv
      {\TIn{(-T)}S}{\TOut TS} \kinds}

    \infrule{\rulenametconvNegOut}{\typeAgainst T \kindp \\ \typeAgainst S \kinds}{\isConv
      {\TOut{(-T)}S}{\TIn TS} \kinds}
    
    \infrule{\rulenametconvNegInv}{\typeAgainst T \kindp}{\isConv {-(-T)}T \kindp}
  \end{mathpar}
  \caption{Type conversion (congruence rules omitted)}
  \label{fig:type-conversion}
\end{figure}

    % \infrule{\rulenametconvDoubleG}{\isType G \kindg}{\isConv{\TDual{(\TDual G)}} G \kindg}
    % 
    % \infrule{\rulenametconvDualDot}{
    % \isType G \kindg \\ \isType S {\kinds}
    % }{
    %   \isConv{\TDual{(G.S)}}{\TDual G . \TDual S} \kinds
    % }
    %   
    %   \infrule{}{ \isType{\forall
    %   (\typec{\alpha}\colon\kindc{\kind}) \typec{T_1}}{\kindc{\kind'}} \\ \kindc{\kind'} = \kinds,\kindg }
    %   { \isConv {\TDual{(\forall
    %   (\typec{\alpha}\colon\kindc{\kind}) \typec{T_1} })} {\forall
    %   (\typec{\alpha}\colon\kindc{\kind}) \typec{(\TDual{T_1})}}  { \kind' }}
    %   
    %   \\
    % \infrule{\rulenametconvDualOutP}{\typeAgainst T \kindp }{ \isConv{\TDual{(!T)}} {?(\TDual{T}) } \kindg}
    % 
    % \infrule{\rulenametconvDualInP}{\typeAgainst T \kindp}{ \isConv{\TDual{(?T)}} {!(\TDual{T})} \kindg}
    % 
    % \infrule{\rulenametconvDualProto}{
    % % \protocols \\ 
    % \typeAgainst{P\ \overline{T}}\kindp }{
    % \isConv{\TDual{(P\ \overline{T}) }}{P\ \overline{T} }{\kindp} }
    % 
    % \infrule}
    % { \typeAgainst{T} {\kinds,\kindg} }
    % {\TDualInner{T} = \TDual{T} }
    % 
    % \infrule{}
    % { \typeAgainst{T} {\kindm,\kindp} }
    % {\TDualInner{T} = \typec{T} }
    %
    % \infrule{\rulenametconvPosOut}{\typeAgainst T \kindp \\ \typeAgainst S \kinds}{\isConv {!{+T}.S}{!T.S} \kinds}

    % \infrule{\rulenametconvPosIn}{\typeAgainst T \kindp \\ \typeAgainst S \kinds}{\isConv {?{+T}.S}{?T.S} \kinds}
    % \\
    % \infrule{\rulenametconvPos}{\typeAgainst T \kindp}{\isConv {+T}T \kindp}
    % \infrule{\rulenametconvReflex}{\typeAgainst T \kind}{\isConv T T \kind}
    %
    % \textcolor{red}{
    %   \infrule{\rulenametconvSemiOut}{
    %     \typeAgainst T \kindp \\ \typeAgainst U \kindp \\ \typeAgainst S \kinds
    %   }{
    %     \isConv{!(T;U).S}{!T.!U.S} \kinds
    %   }
    % }
    %   
    % \textcolor{red}{
    %   \infrule{\rulenametconvSemiIn}{
    %     \typeAgainst T \kindp \\ \typeAgainst U \kindp \\ \typeAgainst S \kinds
    %   }{
    %     \isConv{?(T;U).S}{?T.?U.S} \kinds
    %   }
    % }


%%% Local Variables:
%%% mode: latex
%%% TeX-master: "main"
%%% End:


Types obey a kind-indexed conversion
relation. \Cref{fig:type-conversion} presents the declarative
rules. Conversion is reflexive, symmetric, and transitive
(\rulenametconvReflex, \rulenametconvSymm, \rulenametconvTrans).
%
We define duality as a conversion on session types. The rules entail
that duality changes the direction of transmissions from the outside
(\rulenametconvDualEndW, \rulenametconvDualEndT,
\rulenametconvDualInM, \rulenametconvDualOutM)
and that duality is involutory (\rulenametconvDoubleS). Kind subsumption is lifted to conversion (\rulenametconvSub). The last line
governs the interaction between the reverse operator and the session
type constructors: reverse flips the direction from the inside
(\rulenametconvNegIn, \rulenametconvNegOut). Reverse is
also involutory (\rulenametconvNegInv). We omit the standard congruence rules.


%%% TYPE OPS _ Moved to expression typing

% The following two type operators are used to \emph{polarize} the protocol's invocation.
% \pt{extended the metafunctions ``!'' and ``?'' to $O$ by making them
%   act correctly on guards.}

% \begin{figure}[t!]
\declrel{Materialization and the directional operators}
{$\tosession TT = \typec{T} \quad 
  \polarizedparam - T = \typec{T} \quad
  \polarizedparam + T = \typec{T}$}
%
\begin{align*}
  \tosession {\pneg T} U &= \TIn TU & 
  \polarizedparam - {\pneg T} &= \polarizedparam + T &
  \polarizedparam + {\pneg T} &= \polarizedparam - T
  \\
  \tosession TU &= \typec{\TOut TU} & 
  \polarizedparam - T &= \pneg T &
  \polarizedparam + T &= \typec{T}
\end{align*}
In the second line $\typec T \neq \pneg U$ in all cases
\caption{Auxiliary metafunctions on session types and protocol types}
\label{fig:type-polarized-operator}
\end{figure}


%%% Local Variables:
%%% mode: latex
%%% TeX-master: "main"
%%% End:


% Rather obvious
% \begin{lemma}[Duality is for sessions]
%   \label{lem:duailty-st}
%   If $\typeSynth{\TDual{T}}{\kind}$, then $\typeAgainst{\TDual{T}}{\kinds}$ and
%   $\typeAgainst{T}{\kinds}$.
% \end{lemma}

While type formation is algorithmic, the declarative definition of
type conversion is more perspicuous. For the
upcoming algorithmic expression typing relation we define a normalization
algorithm for types, which is sound and complete with respect to the
declarative rules.

% We start by specifying the syntax of normal forms. Normal forms for well-formed types
% are $\typec{Q}$ as defined by the following grammar:
% %
% \begin{align*}
%   \typec{Q} &\grmeq \typec{R} \grmor \pneg{R}
%   \\
%   \typec{R} &\grmeq
%   % functional types
%   \TUnit \grmor \TFun[]{R}{R} \grmor \TPair{R}{R} \grmor
%   \TForall \alpha \kind {R} \grmor \typec\alpha
%   % session types
%   \grmor \TIn{R}{R} \grmor \TOut{R}{R}
%   \grmor \TEndW \grmor \TEndT  \grmor
%   \TDual{\alpha}
%   % protocol and data types
%   \grmor \TPApp \proto {\overline{Q}}
% \end{align*}
In a type in normal form,  the reverse operator can only occur at most once at the top level (of
a type of kind $\kindp$) and at the top level of the parameters of a protocol type. The
$\dualk$ operator only appears on type variables, at the end of
the spine of a session type. Thus an algorithm to compute normal forms
must push all $\dualk$ operators down the spine of a session type
by applying the conversion rules \rulenametconvDualInM, \rulenametconvDualOutM,
\rulenametconvDualEndW, and \rulenametconvDualEndT from left to
right. If it encounters another $\dualk$ operator along the
spine, it applies rule \rulenametconvDoubleS. We cannot further resolve
$\TDual\alpha$, hence it is a normal form.

The reverse operator only applies to protocol types. So it can only
occur in a message type like $\TIn{(-T)}S$ and in the parameter of a
protocol type. In a message type, normalization must apply rules
\rulenametconvNegIn and \rulenametconvNegOut exhaustively from left to
right, which removes the reverse operator from message types. In the
parameter of a protocol type, normalization applies rule
\rulenametconvNegInv exhaustively. Starting with an odd number of
reverse operators, a single operator remains; starting with an even
number, all reverse operators are removed.


%\begin{figure}[t!]
\declrel{Materialization and the directional operators}
{$\tosession TT = \typec{T} \quad 
  \polarizedparam - T = \typec{T} \quad
  \polarizedparam + T = \typec{T}$}
%
\begin{align*}
  \tosession {\pneg T} U &= \TIn TU & 
  \polarizedparam - {\pneg T} &= \polarizedparam + T &
  \polarizedparam + {\pneg T} &= \polarizedparam - T
  \\
  \tosession TU &= \typec{\TOut TU} & 
  \polarizedparam - T &= \pneg T &
  \polarizedparam + T &= \typec{T}
\end{align*}
In the second line $\typec T \neq \pneg U$ in all cases
\caption{Auxiliary metafunctions on session types and protocol types}
\label{fig:type-polarized-operator}
\end{figure}


%%% Local Variables:
%%% mode: latex
%%% TeX-master: "main"
%%% End:

\begin{figure}[t!]
    \declrel{Type normalization}{$\Normal[+]T=\typec T \quad \Normal[-]T=\typec T$}

  \begin{minipage}{0.49\linewidth}
    \begin{align*}
    % Positive normalisation
    % Functional
    \Normal[+]\TUnit &= \TUnit \\
    \Normal[+]{\TFun[] TU} &= \TFun[]{\Normal[+] T}{\Normal[+] U} \\
    \Normal[+]{\TPair TU} &= \TPair{\Normal[+] T}{\Normal[+] U} \\
    \Normal[+]{\TForall \alpha \kind T} &= \TForall\alpha \kind {\Normal[+] T} \\
    % Functional and session
    \Normal[+]\alpha &= \typec\alpha \\
    % Session
    \Normal[+]{\TIn TS} &= \TFIn{\Normal[+]T}{\Normal[+] S} \\
    \Normal[+]{\TOut TS} &= \TFOut{\Normal[+]T}{\Normal[+]S} \\
    \Normal[+]\TEndW &= \TEndW \\
    \Normal[+]\TEndT &= \TEndT\\
    \end{align*}
  \end{minipage}
  \begin{minipage}{0.49\linewidth}
    \begin{align*}
    \Normal[+]{\TDual T} &= \Normal[-]{T} \\
    % Protocol
    \Normal[+]{\TPApp \proto{\overline T}} &= \typec{\TPApp \proto {\overline{\Normal[+] {T}}}} \\
    \Normal[+]{\pneg T} &= \polarizedparam - {\Normal[+] T}\\
    \Normal[-]{\TDual T} &= \Normal[+]{T} \\
    % Negative noralisation
    \Normal[-]\alpha &= \TDual\alpha \\
    \Normal[-]{\typec{\TIn TS}} &= \TFOut{\Normal[+]T}{\Normal[-]S} \\
    \Normal[-]{\typec{\TOut TS}} &= \TFIn{\Normal[+]T}{\Normal[-]S} \\
    \Normal[-]\TEndW &= \TEndT \\
    \Normal[-]\TEndT &= \TEndW \\
    \end{align*}
  \end{minipage}

  \declrel{Materialization and the directional operators}
  {$\tosession TT = \typec{T} \quad 
    \polarizedparam - T = \typec{T} \quad
    \polarizedparam + T = \typec{T}$}
  %
  \begin{align*}
    \tosession {\pneg T} U &= \TIn TU & 
    \polarizedparam - {\pneg T} &= \polarizedparam + T &
    \polarizedparam + {\pneg T} &= \polarizedparam - T
    \\
    \tosession TU &= \typec{\TOut TU} & 
    \polarizedparam - T &= \pneg T &
    \polarizedparam + T &= \typec{T}
  \end{align*}
  In the second line $\typec T \neq \pneg U$ in all cases.
  \caption{Type normalization and auxiliary metafunctions on session types and protocol types}
  \label{fig:normalisation}
  \label{fig:type-polarized-operator}
\end{figure}

%%% Local Variables:
%%% mode: latex
%%% TeX-master: "main"
%%% End:


Following this intuition normalization is defined in \cref{fig:normalisation} by two mutually
recursive functions $\Normal[+]\_$ and $\Normal[-]\_$ along with three
auxiliary functions. % also  in \cref{fig:type-polarized-operator}.
Positive normalization $\Normal[+]\_$ traverses and reconstructs all
non-session type constructs. It also gets invoked on the type of the
transmitted value for session types, i.e., on $\typec T$ in type
$\TOut TS$.
Normalizing the $\dualk$ operator switches forth and back  between
positive and negative normalization, $\Normal[+]\_$ and $\Normal[-]\_$.

Negative normalization with superscript
$^-$ indicates a pending $\dualk$ type constructor, so that
the function $\Normal[-]\_$ is only used for session
types. Normalization pushes the pending dual constructor along the spine of a session type to
implement rules \rulenametconvDualInM, \rulenametconvDualOutM,
\rulenametconvDualEndW, and \rulenametconvDualEndT. The pending dualization gets reified on
a type variable.

Normalizing a reversed protocol $\TMinus T$ invokes the negative directional
function $\polarizedparam - {T'}$ on $\typec
T$'s normal form to enforce the
conversion rule \rulenametconvDoubleS.

% As an example for the
% normalization of session types, we consider $\Normal{\TIn TS}$. In each case, we invoke $\Normal S$ on
% the continuation session and invoke $\Normal[+]T$ on the transmitted type
% $\typec T$. For sending (receiving), we invoke the positive (negative) directional function
% on the result. Materialization $\tosession TS$ expects a normal form for
% $\typec T$ and ``materializes'' the direction of the session type. It returns an input $\TIn{T'}S$ if the normal form is
% negative and an output $\TOut{T'}S$ if it is positive.
%
Here is an example of a step-by-step computation of the positive normal form of
 $\TDual{(\TIn{(-\TInt)}{\alpha})}$.
%\Cref{fig:example-computation-normal-form} goes through a concrete normalization step by step.
%
% \begin{figure}
  \begin{align*}
    & \Normal[+]{\TDual{(\TIn{(-\TInt)}{\alpha})}} \\
    ={}& \Normal[-]{{\TIn{(-\TInt)}{\alpha}}}  && \text{negative
                                                  normalization
                                                  remembers pending $\dualk$}\\
    ={}& \TFOut{\Normal[+]{-\TInt}}{\Normal[-]\alpha}  &&
                                                          \text{$\Normal[+]\_$
                                                          on the
                                                          payload,
                                                          $\Normal[-]\_$
                                                          on the spine}\\
    ={}& \TFOut{\TFMinus{\Normal[+]{\TInt}}}{\TDual\alpha} &&
                                                              \text{negative
                                                              normalization
                                                              reifies
                                                              $\dualk$
                                                              on type variable}\\
    % ={}& \TFOut{\TFMinus{\TInt}}{\TDual\alpha} \\
    % ={}& \TFOut{-{\TInt}}{\TDual\alpha} \\
    ={}& \tosession{\TMinus\TInt}{\TDual\alpha} &&
                                                   \text{materialization
                                                   applied to normal
                                                   form \dots}\\
    ={}& \TIn{{\TInt}}{\TDual\alpha} && \text{\dots{} fixes the direction}
  \end{align*}
%   \caption{Example: Step-by-step computation of the positive normal form of
%   $\TDual{(\TIn{(-\TInt)}{\alpha})}$}
%   \label{fig:example-computation-normal-form}
% \end{figure}


To test whether a type equivalence goal $\isConv TU\kind$ holds, we
compute
\emph{algorithmic type equivalence}, notation $\isEquiv TU$, which is
defined by comparision of normal forms: $\Normal[+]T = \Normal[+]U$  up to  $\alpha$-equivalence.
% We write $\typec T =_\alpha \typec U$ if types $\typec T$ and $\typec U$ are
% $\alpha$-equivalent.
% 
The following propositions establish soundness and
completeness of the normalization algorithm, as well as its worst case running
time.
% , which is linear in the size of the input type.

% Every type is convertible to its normal form.
%
\begin{theorem}[Soundness]
  \label{thm:soundness-type-equiv}
  Let $\typeSynth T\kind$ and $\typeSynth U\kind$.
  \begin{enumerate}
  \item\label{it:soundness-type-equiv-t} If $\Normal[+] T =_\alpha \Normal[+] U$,
    then $\isConv TU\kind$.
  \item\label{it:soundness-type-equiv-s} If $\kind = \kinds$ and  $\Normal[-] T =_\alpha \Normal[-]
    U$, then $\isConv TU\kinds$.
  \end{enumerate}
\end{theorem}

% If two types are convertible, then their normal forms are equal.
\begin{theorem}[Completeness]
  \label{thm:completeness-type-equiv}
  Let $\isConv TU\kind$.
  \begin{enumerate}
  \item  $\Normal[+]T =_\alpha \Normal[+]U$.
  \item    If $\kind=\kinds$, then $\Normal[-]T =_\alpha \Normal[-]U$.
  \end{enumerate}
\end{theorem}

\begin{theorem}[Time complexity of type equivalence]
  \label{thm:complexity}
  The test for\/ $\isEquiv TU$ runs in $O(|\typec T|+|\typec U|)$.
\end{theorem}


%%% Local Variables:
%%% mode: latex
%%% TeX-master: "main"
%%% End: